% Chapter 1



%----------------------------------------------------------------------------------------

% Define some commands to keep the formatting separated from the content 
%\newcommand{\keyword}[1]{\textbf{#1}}
%\newcommand{\tabhead}[1]{\textbf{#1}}
%\newcommand{\code}[1]{\texttt{#1}}
%\newcommand{\file}[1]{\texttt{\bfseries#1}}
%\newcommand{\option}[1]{\texttt{\itshape#1}}

%----------------------------------------------------------------------------------------

\section{路加福音}
\subsection{天國的十六篇比喻}
路加特别强调的是主到耶路撒冷去的旅程，三卷书比较起来，马太只有两章的篇幅，而马可只用一章篇幅来报导这事。但路加却用了超过十章的篇幅来，这是路加福音最长的一段了（9：51—19：44）

其中的十六篇比喻

1：好撒玛利亚人的比喻

2：再三切求的比喻

3：无知财主的比喻

4：不结实的无花果树的比喻

5：席上客人与主人让位的比喻

6：大筵席的比喻

7：计算花费盖楼的比喻

8：酌量出兵打仗的比喻

9：三重比喻：失羊，失财，浪子（对象是法利赛人与文士）

10：不义管家的比喻

11：财主与拉撒路的故事

12：讲解主人与仆人的故事

13：医治十个长大麻疯的人

14：不义的官的比喻

15：法利赛人与税吏的祷告的比喻

16：仆人与银子的比喻

\subsection{兩個嬰孩(1:5-2:21)}
預言施洗約翰生(1:5-25)、
預言耶穌生(1:26-38)、
兩位母親的見面(1:39-56)、
施洗約翰生(1:57-80)、
耶穌誕生(2:1-21)。

\subsection{兩個見證(2:22-52)}
西面見耶穌(2:22-35)、
女先知亞拿見耶穌(2:36-39)、耶穌的小時候(2:37-52)。


\section{耶穌的講論(按照內容分類順序)}
\subsection{天國的講論}

\subsubsection{路加福音關於天國的十六篇比喻}
路加特别强调的是主到耶路撒冷去的旅程，三卷书比较起来，马太只有两章的篇幅，而马可只用一章篇幅来报导这事。但路加却用了超过十章的篇幅来，这是路加福音最长的一段了（9：51—19：44）

其中的十六篇比喻

1：好撒玛利亚人的比喻

2：再三切求的比喻

3：无知财主的比喻

4：不结实的无花果树的比喻

5：席上客人与主人让位的比喻

6：大筵席的比喻

7：计算花费盖楼的比喻

8：酌量出兵打仗的比喻

9：三重比喻：失羊，失财，浪子（对象是法利赛人与文士）

10：不义管家的比喻

11：财主与拉撒路的故事

12：讲解主人与仆人的故事

13：医治十个长大麻疯的人

14：不义的官的比喻

15：法利赛人与税吏的祷告的比喻

16：仆人与银子的比喻


\subsubsection{天國的君王}
 \begin{table}[H]
  \centering
 % \caption{Add caption}
    \begin{tabular}{|p{3.93em}|p{11.53em}|c|p{4.045em}|p{4.045em}|p{4.045em}|p{4.045em}|}
    \toprule
    \multicolumn{2}{|p{16.5em}|}{事 件} & 地 點 & 馬太福音 & 馬可福音 & 路加福音 & 約翰福音 \\
    \midrule
    \multicolumn{2}{|p{13.5em}|}{耶穌與尼哥底母談道（逾越節）} & 耶路撒冷 &       &       &       & 3:1-21 \\
  \midrule   
 \multicolumn{2}{|p{11.5em}|}{耶穌與婦人傳道} & 敘加 &       &       &       & 4:3下-38 \\
 \midrule
 \multicolumn{2}{|p{11.5em}|}{耶穌在馬太家坐席} & \multirow{1}{*}{迦百農} & 9:10-13 &2:15-17& 5:29-32 &  \\
\midrule
 \multicolumn{2}{|p{16.5em}|}{耶穌論門徒禁食} &       & 9:14-17 & 2:18-22 & 5:33-39 &  \\
 \midrule
 
    \multicolumn{1}{|l|}{\multirow{4}{*}{\parbox{0.08\textwidth}{耶穌講論父與子的關係（逾越節）}}} & \multicolumn{1}{p{12.715em}|}{a.子照著父所作的事去行} & \multirow{4}[8]{*}{耶路撒冷} & \multicolumn{1}{r|}{} & \multicolumn{1}{r|}{} &       & \multicolumn{1}{p{5.93em}|}{5:17-21 } \\
\cmidrule{2-2}\cmidrule{4-7}    \multicolumn{1}{|l|}{} & \multicolumn{1}{p{12.715em}|}{b.父把審判的事交給子} & \multicolumn{1}{c|}{} & \multicolumn{1}{r|}{} & \multicolumn{1}{r|}{} &       & \multicolumn{1}{p{5.93em}|}{5:22-30 } \\
\cmidrule{2-2}\cmidrule{4-7}    \multicolumn{1}{|l|}{} & \multicolumn{1}{p{12.715em}|}{c.父爲子作見證} & \multicolumn{1}{c|}{} & \multicolumn{1}{r|}{} & \multicolumn{1}{r|}{} &       & \multicolumn{1}{p{5.93em}|}{5:31-42} \\
\cmidrule{2-2}\cmidrule{4-7}    \multicolumn{1}{|l|}{} & \multicolumn{1}{p{12.715em}|}{d.子奉父的命而來} & \multicolumn{1}{c|}{} & \multicolumn{1}{r|}{} & \multicolumn{1}{r|}{} &       & \multicolumn{1}{p{5.93em}|}{5:43-47} \\
    \midrule

    \bottomrule
    \end{tabular}%
  \label{tab:addlabel}%
    \caption{到第二個逾越節之前的講論}
\end{table}%

\subsubsection{天國的子民}

%\subsection{門徒成熟講論(按照內容分類順序)}

\subsection{耶穌關於預言的講論}
\subsubsection{預言彼得不忍主}
\subsubsection{預言猶大出賣主}
 \subsubsection{預言受死復活（四次）}
 \subsubsection{復活後去加利利}
 
 %\subsection{耶穌關於天國的講論}

\subsection{耶穌的責備(按照內容分類順序)}
 \subsubsection{責備法利賽人}
\subsubsection{責備門徒}
 
 







